\chapter{Common Git Operations}
\textbf{git init}

初始化仓库命令,新建文件夹的第一步操作

\section{Configuration}
\textbf{更换Git Bash默认文本编辑器}
\begin{myminted}{bash}
    git config --global core.editor Code
\end{myminted}

\textbf{user and email configuration}
\begin{myminted}{bash}
    git config --global user.name "Your Name"
    git config --global user.email "your.email@example.com"
\end{myminted}

\textbf{configuration checkout}
\begin{myminted}{bash}
    git config --list
\end{myminted}

\textbf{SSH configuration}

\section{Basic Operations}
\begin{myminted}{bash}
    git add . #添加所有文件
    git commit -m "commit statement" #提交更改

    # 创建.gitignore文件
    code .gitignore #code 为git bash 默认的文本编辑器
\end{myminted}


\section{Remote Operations}
\begin{myminted}{bash}
    git remote add <remote-name> <remote-url> #添加远程仓库，<remote-name>是自定义名称，<remote-url>是远程仓库的URL

    git remote -v #查看远程仓库
    git push -u <remote-name> <branch-name> #推送到远程

    git push #再次提交后可以省略,注意要完成git add和commit操作
\end{myminted}

\textbf{pull remote repository}

\section{Branch Operations}
\begin{myminted}{bash}
    git branch  #查看本地分支
    git branch -a  #查看所有分支(含远程)
\end{myminted}

\section{Undo \& Revert}
\begin{myminted}{bash}
    git reset HEAD <file> #撤销暂存区文件
\end{myminted}

\section{Differences}
\begin{myminted}{bash}
    git diff #查看工作区与暂存区的差异
\end{myminted}